\begin{rubric}{Employment}

\entry*[2026--now]
  \role{Staff AI Engineer} --- \textbf{iFood Tech} (Remote, Brazil)
  \begin{itemize}
  \item Building \emph{N1ago}, a post-sales AI agent handling consumer-credit interactions at scale:
    LLM-powered conversational flows behind a pluggable provider layer, with safe defaults when
    customer data is missing. Stack: Python (FastAPI), React, PostgreSQL, Redis, cloud services
    and CI/CD for scalable deployment.
  \item Fine-tuned a GPT model to classify N1ago's responses, replacing brittle heuristics with a
    task-specific classifier.
  \item Built the evaluation stack --- a golden dataset that regression-tests LLM responses across
    prompt and model changes, making quality shifts measurable before release.
  \item Instrumented LLM trace observability (Langfuse) and APM (Datadog), making prompt
    regressions and latency issues debuggable in production.
  \end{itemize}

\entry*[2021--2026]
  \textbf{Zup Innovation} (Remote, Brazil) --- \role{Staff AI Engineer} (2023--2026),
  \role{Head of Research} (2021--2023)
  \begin{itemize}
  \item Led a team of five researchers; defined the research agenda and turned it into functional
    proofs-of-concept validated with engineering and product teams. Early adopter of LLMs inside
    the company (alongside encoder-only transformers and classical ML), before they went mainstream.
  \item Built and maintained a Retrieval-Augmented Generation system in Python (LangChain, FastAPI,
    Flask, Docker), with vector databases (\tech{pgvector}, \tech{ChromaDB}) storing
    high-dimensional embeddings for retrieval at scale.
  \item Created the \textbf{benchmark framework} used to evaluate StackSpot AI agents --- the
    reference harness for measuring agent quality across releases.
  \item Engineered agentic workflows (LangChain, LangGraph, StackSpot AI) over open and closed
    LLMs/SLMs, with API orchestration, cloud services and CI/CD for scalable deployment.
  \item Prototyped end-to-end AI use cases --- ingestion, embedding generation, retrieval,
    downstream evaluation --- validating feasibility before hand-off to product.
  \end{itemize}

\entry*[2017--now]
  \role{Assistant Professor} --- \textbf{Federal University of Pará (UFPA)} (Belém, Brazil)
  \begin{itemize}
  \item Published 100+ research papers, most in highly competitive venues (ICSE, TSE, EMSE, CACM),
    with artifacts and replication packages released alongside them.
  \item Supervised 40+ students (4 PhD, 10 MSc, 10+ capstone, 10+ undergraduate research);
    raised 800k+ BRL as PI and 2M+ BRL as co-PI for R\&D on AI for software engineering.
  \item Editor-in-Chief of JSERD (2020--2022), leading a board of around 30 associate editors;
    CNPq Research Productivity Fellow since 2020 (renewed in 2023 and 2026).
  \end{itemize}

\end{rubric}
